\begin{abstract}
Retrieval-augmented generation (RAG) and graph-based retrieval (GraphRAG)
are the dominant paradigms for grounding LLM responses in structured
knowledge. Both optimize for recall while treating token cost as a
secondary concern. We introduce Compact Knowledge Graphs (CKG)---pre-structured
DAG representations with explicit concept taxonomy and pipe-delimited
dependency encoding---and present a two-track benchmark comparing CKG,
RAG, and GraphRAG across educational and commercial domains.

The current analyzed artifacts include 47 CKG domains (8,121 queries),
44 RAG domains (7,414 queries), and 20 GraphRAG domains (2,683 queries).
CKG achieves macro-average F1 of \textbf{0.4926} versus \textbf{0.1223}
for RAG and \textbf{0.1200} for GraphRAG, at \textbf{11$\times$ fewer
tokens per query} (264 vs.\ 2,986 vs.\ 3,450). CKG F1 increases with
hop depth to 0.772 while RAG remains substantially lower. The compound
Reasoning Density Score (RDS) advantage is \textbf{45$\times$}. T1
entity lookup (CKG 0.222, RAG 0.093)
serves as a designed negative control.

Track 2 tests whether the architecture generalizes beyond hand-curated
educational data. We build a complete GLP-1/Obesity pharmacology CKG
programmatically from the ClinicalTrials.gov API with no expert
annotation (170 queries). The pipeline-generated CKG reaches macro F1
\textbf{0.5298}---exceeding the current CKG average by 7.5\%
and preserving a \textbf{28$\times$} RDS advantage over RAG. Automated
construction does not degrade retrieval quality; it matches or improves
it. Together the two tracks establish CKG as a distinct architecture
whose structural advantage is domain-agnostic and does not depend on
manual curation.

Benchmark, dataset, and evaluation harness are released at
\url{https://github.com/Yarmoluk/ckg-benchmark}
under CC BY 4.0\,/\,MIT.
\end{abstract}
